\section{符号}
有些作者在曲线上添加符号以向读者提供更多信息。以下是几种可能的用法（\TIKZ\ 库 \tkzname{arrows}）：

\begin{tikzpicture}
\draw[thick,(-)](0,0)--(2,2);
\draw[thick,o-o](2,0)--(4,2);
\draw[thick,)-(](4,0)--(6,2);
\draw[thick,*-*](6,0)--(8,2);
\end{tikzpicture}

\newcommand{\cred}[1]{{\color{red}#1}}
\newcommand{\cgreen}[1]{{\color{green!50!black}#1}}
\newcommand{\cblue}[1]{{\color{blue}#1}}

以下示例由 \tkzname{Simon Schläpfer} 提供：

我们要绘制
\[
y=\left\{\begin{array}{ll}
    \cred{8-1.5x}&,\text{当 }x<2\\
    \cblue{4}&,\text{当 }2 \leq x \leq 3\\
    \cgreen{2x-4}&,\text{当 } x>3
  \end{array}
\right.
\]

曲线不能裁剪到窗口内（\tkzname{clip=false}），否则端点处的符号会被截断。

\begin{center}
\begin{tkzexample}[vbox]
\begin{tikzpicture}
  \tkzFnInit[xmin=-1,xmax=6,ymin=0,ymax=10]
  \tkzFnGrid[color=gray]
  \tkzFnAxes
  \tkzFnDraw[{-[},clip=false,color=red,domain=-1:2,samples=2]{8-1.5*x}
  \tkzFnDraw[{[-]},clip=false,color=blue,domain=2:3,samples=2]{4}
  \tkzFnDraw[{]-},clip=false,color=green!50!black,domain=3:6,samples=2]{2*x-4}
\end{tikzpicture}
\end{tkzexample}
\end{center}

\endinput
